\chapter{Factors, Matrices, and Arrays}
\label{chap:factors-matrices-arrays}

\chapterguide
  {You should be comfortable with vectors, lists, indexing, and arithmetic.}
  {construct and inspect factors; access factor elements and levels; construct matrices by row or column; name and index matrix dimensions; apply element-wise matrix arithmetic; create named multidimensional arrays; access slices and elements; and use \texttt{apply()} as shown in the handout.}

\section{Factors}

Lab 4a creates a character vector containing compass regions and then converts it to a factor:

\begin{lstlisting}[style=dsr]
data <- c(
  "East", "West", "East", "North", "North",
  "East", "West", "West", "West", "East", "North"
)

print(data)
print(is.factor(data))

factor_data <- factor(data)
print(factor_data)
print(is.factor(factor_data))
length(factor_data)
\end{lstlisting}

The before/after \texttt{is.factor()} checks are the key observation: the original vector is not a factor, while the result of \texttt{factor()} is.

\subsection{Accessing and modifying factor values}

\begin{lstlisting}[style=dsr]
data <- factor(c(
  "East", "West", "East", "North", "North",
  "East", "West", "West", "West", "East", "North"
))

data[3]
data[3] <- "NorthWest"
\end{lstlisting}

The modification attempt is intentionally useful: it tests what happens when a value not already represented by the factor's levels is assigned.

\subsection{Changing level order}

\begin{lstlisting}[style=dsr]
data <- c(
  "East", "West", "East", "North", "North",
  "East", "West", "West", "West", "East", "North"
)

factor_data <- factor(data)
new_order_data <- factor(
  factor_data,
  levels = c("East", "West", "North")
)
print(new_order_data)
\end{lstlisting}

This exercise separates the observed values from the order in which the levels are defined.

\subsection{Generating factor levels}

\begin{lstlisting}[style=dsr]
v <- gl(3, 4, labels = c("Tampa", "Seattle", "Boston"))
print(v)
\end{lstlisting}

The handout uses \texttt{gl()} to generate repeated factor levels programmatically.

\section{Matrices}

A matrix organizes values into rows and columns. Lab 4a contrasts row-wise and column-wise filling:

\begin{lstlisting}[style=dsr]
M <- matrix(c(3:14), nrow = 4, byrow = TRUE)
print(M)

N <- matrix(c(3:14), nrow = 4, byrow = FALSE)
print(N)
\end{lstlisting}

\subsection{Naming rows and columns}

\begin{lstlisting}[style=dsr]
rownames <- c("row1", "row2", "row3", "row4")
colnames <- c("col1", "col2", "col3")

P <- matrix(
  c(3:14),
  nrow = 4,
  byrow = TRUE,
  dimnames = list(rownames, colnames)
)
print(P)
\end{lstlisting}

Dimension names make later inspection easier because row and column labels appear with the values.

\subsection{Matrix indexing}

\begin{lstlisting}[style=dsr]
print(P[1, 3])  # row 1, column 3
print(P[4, 2])  # row 4, column 2
print(P[2, ])   # all columns from row 2
print(P[, 3])   # all rows from column 3
\end{lstlisting}

The index order in the supplied code is \texttt{[row, column]}.

\section{Matrix arithmetic in the handout}

Two 2-by-3 matrices are added, subtracted, multiplied, and divided using ordinary arithmetic operators:

\begin{lstlisting}[style=dsr]
matrix1 <- matrix(c(3, 9, -1, 4, 2, 6), nrow = 2)
matrix2 <- matrix(c(5, 2, 0, 9, 3, 4), nrow = 2)

print(matrix1 + matrix2)
print(matrix1 - matrix2)
print(matrix1 * matrix2)
print(matrix1 / matrix2)
\end{lstlisting}

\begin{keyidea}
The multiplication and division examples in Lab 4a use \texttt{*} and \texttt{/}. In these notes, ``matrix multiplication'' in the Lab 4 practical is kept consistent with that supplied element-wise operator usage rather than silently replacing it with a different operation.
\end{keyidea}

\section{Arrays}

Lab 4a extends the row/column idea into a third dimension.

\subsection{Constructing an array}

\begin{lstlisting}[style=dsr]
vector1 <- c(5, 9, 3)
vector2 <- c(10, 11, 12, 13, 14, 15)

result <- array(
  c(vector1, vector2),
  dim = c(3, 3, 2)
)
print(result)
\end{lstlisting}

The dimensions request 3 rows, 3 columns, and 2 matrix-like layers.

\subsection{Naming array dimensions}

\begin{lstlisting}[style=dsr]
column.names <- c("COL1", "COL2", "COL3")
row.names <- c("ROW1", "ROW2", "ROW3")
matrix.names <- c("Matrix1", "Matrix2")

result <- array(
  c(vector1, vector2),
  dim = c(3, 3, 2),
  dimnames = list(row.names, column.names, matrix.names)
)
print(result)
\end{lstlisting}

\subsection{Accessing array elements and slices}

\begin{lstlisting}[style=dsr]
print(result[3, , 2])  # third row of second matrix
print(result[1, 3, 1]) # row 1, column 3, first matrix
print(result[, , 2])   # entire second matrix
\end{lstlisting}

The indexing pattern generalizes matrix indexing: \texttt{[row, column, layer]}.

\section{Array manipulation and \texttt{apply()}}

The handout extracts 2D matrices from arrays and adds them. It also demonstrates \texttt{apply()} across rows:

\begin{lstlisting}[style=dsr]
new.array <- array(
  c(vector1, vector2),
  dim = c(3, 3, 2)
)

result <- apply(new.array, c(1), sum)
print(result)
\end{lstlisting}

Here the margin argument \texttt{c(1)} tells the supplied example to calculate sums associated with the row dimension across the array.

\begin{pitfall}
In Lab 4a section 3.4, \texttt{vector3} and \texttt{vector4} are created, but the printed code then constructs \texttt{array2} from \texttt{vector1} and \texttt{vector2} again. These notes preserve the handout as a source caution rather than silently assuming which vectors were intended.
\end{pitfall}

\begin{quickcheck}
\begin{enumerate}
  \item Which function converts the region vector into a factor?
  \item What does \texttt{byrow = TRUE} change when constructing a matrix?
  \item In \texttt{P[4,2]}, which coordinate is the row and which is the column?
  \item How many dimensions are specified by \texttt{dim = c(3,3,2)}?
  \item What does \texttt{result[,,2]} select from a three-dimensional array?
\end{enumerate}
\end{quickcheck}

\begin{chapterrecap}
\begin{itemize}
  \item Factors represent categorical values and expose levels that can be ordered explicitly.
  \item Matrices are built with row/column dimensions and can carry dimension names.
  \item The handout's matrix arithmetic uses ordinary operators element by element.
  \item Arrays generalize matrix indexing to additional dimensions.
  \item \texttt{apply()} is introduced as a way to run a function across a selected array margin.
\end{itemize}
\end{chapterrecap}
